%-------------------------------------------------------------------------------------
\section{Lesson 2, 28/09/2026}
\subsection{Again on SU(2)}
The matrices of SU(2) can be seen as matrices in the form
$$g=
\begin{pmatrix}
    x_0+ix_3&x_2+ix_1\\
    -x_2+ix_1&x_0-ix_3
\end{pmatrix}\text{, with }
x_0^2+x_1^2+x_2^2+x_3^2 = \det(g) = 1
$$
It can be also represented by bigger matrices, but some of them are trivial, like
$$g=
\begin{pmatrix}
    x_0+ix_3&x_2+ix_1&0\\
    -x_2+ix_1&x_0-ix_3&0\\
    0&0&1
\end{pmatrix}\text{, with }
x_0^2+x_1^2+x_2^2+x_3^2 = \det(g) = 1
$$
which is the same as before, but with a bunch of zeros.
\subsection{Generators}
The generators are defined as the "directions" in which matrices can "move away"
from the identity matrix: it is clear that
$$g = x_0\mathbb{I} + i x_i\sigma^i$$
where $\sigma^i$ are the Pauli matrices; since this is a continuous group, the
infinitesimal movement in the direction $\sigma^1$ can be written as
$$g=\mathbb{I} + i\varepsilon\sigma^1 + O(\varepsilon^2).$$

Consider the set of generators of a continuous (Lie) group $\left\{T^a\right\}$: the
commutator between them is
$$[T^a,T^b] = T^aT^b-T^bT^a\underset{Lie}= if\indices{^a^b_c}T^c$$
in the case of $SU(2)$, the structure constants turn out to be
$f\indices{^a^b_c} = \varepsilon\indices{^a^b_c}$, which is the Levi-Civita symbol.

The algebra of the generators of $SO(3)$ is exactly the same.

\subsection{Lorentz-Poincaré Group}
This is the group of transformations of Special Relativity: using 4-vectors, a kind of
transformation that obeys the rules of SR is a boost along the x-axis
$$B_1=\begin{pmatrix}
\gamma&-\beta\gamma&0&0\\
-\beta\gamma&\gamma&0&0\\
0&0&1&0\\
0&0&0&1
\end{pmatrix}
$$
this is a kind of rotation in a Minkowski space; regular spacial rotations also obey
SR
$$\begin{pmatrix}
1&0&0&0\\
0&&&\\
0&&R&\\
0&&&
\end{pmatrix}
$$
where $R\in SO(3)$.
There are two other transformations, the parity and time-inversion:
$$\begin{pmatrix}
1&0&0&0\\
0&-1&0&0\\
0&0&-1&0\\
0&0&0&-1
\end{pmatrix}\qquad
\begin{pmatrix}
-1&0&0&0\\
0&1&0&0\\
0&0&1&0\\
0&0&0&1
\end{pmatrix}
$$

This group has 6 generators: 3 rotations from $SO(3)$ and 3 boosts (along the 3 axes);
what are the generators for the boosts?

From their definitions,
$$\gamma^2 - (\gamma\beta)^2 = 1$$
so they can be parametrized like
$$\begin{dcases*}
\gamma = \cosh(\eta) = 1 +O(\eta^2)\\
\beta\gamma = \sinh(\eta) = \eta + O(\eta^2)
\end{dcases*}
$$
where $\eta$ is called rapidity.

$$B_1(\eta) = \mathbb{I} + i\eta\begin{pmatrix}
0&i&0&0\\
i&0&0&0\\
0&0&0&0\\
0&0&0&0
\end{pmatrix} + O(\eta^2)
$$

So, the generators of rotations are $\{J_i\}$, with the aforementioned algebra.

Similarly, the generators of the boosts $\{K_i\}$ have commutation relations
$$[K^a,K^b]=-i\varepsilon\indices{^a^b_c}J^c$$
Trying to mix the generators yields
$$[J^a,K^b]=i\varepsilon\indices{^a^b_c}K^c.$$
This can be written in a cleaner way, defining
\begin{gather*}
N_a = \frac{J_a+iK_a}{2}\\
M_a = \frac{J_a+-iK_a}{2},
\end{gather*}
it can be proven that they satisfy 2 closed $SU(2)$ algebras:
\begin{gather*}
    [N^a,N^b]=i\varepsilon\indices{^a^b_c}N^c\\
    [M^a,M^b]=i\varepsilon\indices{^a^b_c}M^c\\
    [N_a,M_b]=0
\end{gather*}
Since $SU(2)$ is the angular momentum group, an "azimuthal number" can be defined
for the two algebras, $p,q$: depending on the values of these eigenvalues the field
satisfying the Lorentz invariance has different properties.
\begin{table}[ht!]
    \centering
\begin{tabular}{c|c|c}
    $p$&$q$&Name\\ \hline
    $0$&$0$&Scalar \\
    $1/2$&$0$&Spin 1/2 right\\
    $0$ & $1/2$&Spin 1/2 left\\
    \end{tabular}
\end{table}\\
The addition of the Poincaré group are spacetime translations, which are related to
the 4-momentum $P^\mu$
\subsection{Internal degrees of freedom}
Some particles are invariant under transformations of "internal" degrees of freedom,
like colour charge in quarks: if these transformations are local, they are called
\textbf{gauge transformations}; locally it is always possible to redefine what the
colour "blue" is for a quark.

Why is it possible to assume local symmetries? Because it makes the mediator fields
come out of the equations (what?).

In the standard model, these internal degrees of freedom are
$$\underbrace{SU(3)}_\text{QCD} \times \underbrace{SU(2) \times U(1)}_
\text{Electroweak}$$

\subsection{Other definitions on groups}
\textbf{Def.} A group $(AG4,\cdot)$ is Abelian if and only if
$$\forall g_1,g_2\in AG4,\;g_1\cdot g_2 = g_2\cdot g_1$$
\textbf{Observation 1} The definition of group implies the uniqueness of the
identity element.\\
\textbf{Proof (by contradiction)} Consider 2 identity elements
$$e\neq f:$$
then their product must be, at the same time
$$\begin{dcases*}
  e\cdot f = f\\
  e\cdot f = e
  \end{dcases*}.$$
This contradicts the assumption $e\neq f \qed$ %not working :(

\textbf{Observation 2} The inverse of every element is unique. The proof is similar
to the previous one

\textbf{Def.} A magma is a set with an internal binary operation.\\
\textbf{Def.} A magma $(M,\cdot)$ is called a semi-group if and only if the binary 
operation is associative\\
\textbf{Def.} A semi-group with an identity element is called a monoid.\\
$$\text{Magma}\supset\text{Semi-group}\supset\text{Monoid}\supset\text{Group}\supset
\text{Abelian Group}$$

The set of a group can be
\begin{itemize}
    \item Finite;
    \item infinite and discrete;
    \item infinite and continuous.
\end{itemize}
The smallest non-trivial group is called $\mathbb{Z}_2$:
$$\mathbb{Z}_2 = \{e,a\}$$
and is called cyclic group of order $2$; in order to specify the algebraic structure,
all the possible products must be determined: a useful tool is a Caylem table
\begin{table}[ht!]
    \centering
    \begin{tabular}{c|cc}
		&e	&a\\ \hline
e	&$e \cdot e$&$e\cdot a$\\
	a	&$a\cdot e$&$a\cdot a$\\
    \end{tabular} 
\end{table}
In this case, the only element that has to be explicitly specified is $a\cdot a$.
This must be $a\cdot a=e$, because supposing $a\cdot a =a$ implies
\begin{gather*}
a\cdot a = a\\
(a\cdot a) \cdot a^{-1} = a\cdot a^{-1}\\
a\cdot (a\cdot a^{-1}) = e\\
a\cdot e = e\\
a = e
\end{gather*}
which is absurd.

This group can be realized in many ways (see lesson 1) and also with
$$(\{\text{False},\text{True}\},\text{XOR})$$

\textbf{Def.} Given two groups $(G_1,\cdot)$, $(G_2,\circ)$, their
\textbf{direct group product} is the group $G_1\times G_2$ made from the set of
ordered pairs
$$(a_1;a_2)\text{, with }a_i\in G_i,$$
and the internal operation $\Box = (\cdot\,;\circ)$ that works as follows
$$(a_1;a_2)\Box(b_1;b_2) := (a_1\cdot b_1;a_2\circ b_2)$$
\subsection{Examples of groups}
\textbf{Integers}:
The set $\mathbb{Z}$ is an Abelian group with the regular sum $+$, identity element
$0$ and inverse $n^{-1}:=-n$.\\

\textbf{Reals}:
The set of real numbers is similarly an Abelian group with the regular sum, but to
construct a group with the multiplication the number $0$ must be taken out:
$(\mathbb{R}\setminus0,\cdot)$ is an Abelian group, with $1$ as the identity and
$a^{-1}=\frac{1}{a}$.\\

$\mathbf{n\times m}$\textbf{ matrices over $\mathbb{R}$}: These sets are groups with
the matrix sum, as they are equivalent to $n\times m$ copies of $(\mathbb{R},+)$.

\textbf{Invertible matrices}: These sets are groups with the matrix multiplication
and the inverse is the matrix inverse.
